%%% Dokumentation der Nutzung von KI-Tools (Kurzfassung aus Textbausteinen)
%%% Zutreffende Textbausteine verwenden, Unzutreffendes bitte löschen.
%%% Weitere Arten der Nutzung, die auf dieser Liste nicht erscheinen, sind
%%% unbedingt vorab mit Ihrer Betreuungsperson zu klären.

\begin{aiusage}
	%\aiusednone                          % keinerlei Hilfsmittel der generativen KI benutzt
	%\aiusedfor{ideas}                    % um erste Ideen für Forschungsfragen und Theorien zu sammeln
	\aiusedfor{literature-search}         % um Literatur zu finden
	%\aiusedfor{literature-understanding} % um Literatur besser nachvollziehen zu können
	\aiusedfor{proofreading}              % um Korrektur zu lesen und den sprachlichen Stil zu verbessern
	\aiusedforother{um Programmcode für Training, Implementierung und Auswertung zu erstellen und zu überarbeiten (Claude),}
	\aiusedforother{um Textabschnitte zu entwerfen, die ich anschließend geprüft und überarbeitet habe, insbesondere Kurzfassung, Abstract, Teile von Kapitel~\ref{sec:analysen} und das Kapitel~\ref{sec:verwandte_arbeiten} (Claude),}
	\aiusedforother{um die Referenzannotation der Testvideos zu erstellen, die ich anschließend selbst geprüft habe (Claude),}
	\aiusedforother{um Trainingsbilder sowie die Testvideos zu erzeugen (Gemini, Veo),}
	\aiusedforother{um die Segmentierungsmasken vorzuannotieren (Roboflow),}
	\aiusedforother{um die englische Fassung der Arbeit zu übersetzen (Claude).}
\end{aiusage}

%%% Alternativ: ausführliche Dokumentation (Langfassung), falls vom Studiengang verlangt.
%%% Dazu die Kurzfassung oben löschen (nur eine der beiden Fassungen ist vorgesehen) und den
%%% folgenden Block einkommentieren. Zutreffende Aktivitäten einkommentieren, Unzutreffendes löschen:
%%% \aiactivity{<Aktivität>}{<KI-Tools>}{<Beschreibung der Nutzung>}
%%% Als KI-Tools die Namen der verwendeten Werkzeuge angeben (z. B. ChatGPT, DeepL Write, Microsoft 365
%%% Copilot, Perplexity, GitHub Copilot), als Beschreibung die Art der Nutzung und die betroffenen
%%% Abschnitte der Arbeit. Die Tabelle ersetzt nicht eine gegebenenfalls nötige nähere Beschreibung
%%% der Nutzung im Text der Arbeit. \aiconfirm{...} bestätigt die beiden Aussagen zum
%%% verantwortungsvollen Umgang mit KI-Tools.
%%% Nutzungen, die keiner Aktivität entsprechen: pro Nutzung ein \aifurtheruse{...} (mehrere
%%% erscheinen als Aufzählung); ohne \aifurtheruse steht dort "Keine".
%%% Die Tabelle enthält standardmäßig alle Aktivitäten (wie die Vorlage der FH OÖ),
%%% mit [compact] nur jene mit Einträgen.

%\begin{aiusagedetails}             % Tabelle mit allen Aktivitäten
%%\begin{aiusagedetails}[compact]   % alternativ: nur Aktivitäten mit Einträgen
%	\aiconfirm{accuracy}            % Ergebnisse der KI-Tools überprüft und ggf. korrigiert
%	\aiconfirm{responsibility}      % Verantwortung liegt bei den Autor*innen
%	%\aiactivity{ideas}{ChatGPT}{Sammeln erster Ideen für Forschungsfragen (Kapitel~1)}
%	\aiactivity{literature}{Perplexity}{Suche nach Literatur zu ... (Kapitel~2)}
%	%\aiactivity{methodology}{ChatGPT}{Vergleich möglicher Evaluierungsmethoden (Kapitel~3)}
%	%\aiactivity{coding}{GitHub Copilot}{Generierung von Testfunktionen (Kapitel~4)}
%	%\aiactivity{data}{ChatGPT}{Erstellung von Skripten zur statistischen Auswertung (Abschnitt~5.1)}
%	%\aiactivity{interpretation}{ChatGPT}{Diskussion möglicher Erklärungen der Ergebnisse (Abschnitt~5.3)}
%	%\aiactivity{structure}{ChatGPT}{Vorschläge zur Gliederung von Kapitel~2}
%	%\aiactivity{generation}{ChatGPT}{Erster Entwurf der Kurzfassung, anschließend selbst überarbeitet}
%	%\aiactivity{translation}{DeepL Translator}{Übersetzung von Zitaten aus dem Englischen (Kapitel~2)}
%	\aiactivity{editing}{DeepL Write}{Korrekturlesen der Kapitel~1--5}
%	%\aiactivity{references}{ChatGPT}{Formatierung einzelner Literaturangaben}
%	%\aifurtheruse{Erzeugung von Testdaten für die Evaluierung (Abschnitt~5.2).}
%	%\aifurtheruse{Erstellung von Diagrammen aus den Messdaten (Abschnitt~5.3).}
%\end{aiusagedetails}
